\FloatBarrier
\newpage
\section{数据的描述}
\label{sec:prob-data-description}

先研究一组已经取得的数据。图表保留形态，数字特征概括位置与分散程度；选择怎样概括，需要联系所研究的变量和问题。本节使用教学数据
\begin{equation}\label{eq:prob-description-data}
 8,\ 9,\ 9,\ 10,\ 10,\ 11,\ 12,\ 15
\end{equation}
贯穿各种表示，单位暂不指定。

\subsection{分布图表}

取值类别较少时，可以用频数表记录每类个数，再除以总数得到频率。数值连续或类别很多时，可以按区间分组。若某组频率为\(f_i\)、组距为\(d_i\)，频率直方图的高度取\(f_i/d_i\)，使每个矩形的面积恰为该组频率。

例如将式\eqref{eq:prob-description-data}分入\([8,10)\)、\([10,12)\)、\([12,16)\)，频数依次为3、3、2。第三组虽有较大组距，其矩形面积仍为\(2/8\)。只比较柱高而忽略组距，会改变频率的解释。

\begin{figure}[!ht]
\centering
\includegraphics[width=.39\linewidth]{probability-histogram}\hfill
\includegraphics[width=.56\linewidth]{probability-histogram-exam}
\caption{左图标注各组频率；右图为零件直径数据，柱顶标注频率与组距之比。}
\label{fig:prob-histograms}
\end{figure}

分组会损失组内位置。需要精确计算时应保留原始数据；只有分组频数时，可以用组中值作近似，并将结果按近似量解释。两列成对数值则可用散点图保留对应关系，回归一节再详细研究。

\subsection{位置特征}

数据\(x_1,\ldots,x_n\)的{\heiti 样本均值}为\(\bar x=n^{-1}\sum_ix_i\)。将数据升序排列为\(x_{(1)}\le\cdots\le x_{(n)}\)，奇数个数据的中位数为中间一项，偶数个则取中间两项的平均；出现次数最多的值称为{\heiti 众数}，可以不唯一。

本书采用以下有限样本百分位数约定：对\(0<p<1\)，若\(np\)不是整数，第\(100p\)百分位数取\(x_{(\lceil np\rceil)}\)；若\(np=k\)为整数，则取\((x_{(k)}+x_{(k+1)})/2\)。第25、50、75百分位数分别记为\(Q_1,Q_2,Q_3\)，其中\(Q_2\)就是中位数。

本节数据的均值为10.5，中位数为10，众数为9和10；\(Q_1=9\)、\(Q_3=11.5\)。以\(Q_1,Q_3\)为箱体两端，在中位数处画线，再标记两端位置，就得到箱线图。本节采用两须延伸到最小值、最大值的简单画法。

\begin{figure}[!ht]
\centering
\includegraphics[width=.64\linewidth]{probability-boxplot}
\caption{同一组数据的箱线图，依次标出最小值、三个四分位数和最大值。}
\end{figure}

均值使\(\sum_i(x_i-c)^2\)最小，这由平方误差分解得到。中位数使\(\sum_i|x_i-c|\)最小：把最小、最大项依次配对，每对的绝对偏差和在\(c\)位于两项之间时最小。偶数时，取中间两项的平均；奇数时，取未配对的中间一项，使其偏差为零。两种选择均能同时使各部分最小。

\subsection{离散程度}

定义描述这组数据的{\heiti 样本方差}与标准差为
\begin{equation}\label{eq:prob-descriptive-variance}
 s_n^2=\frac1n\sum_i(x_i-\bar x)^2
 =\frac1n\sum_ix_i^2-\bar x^2,\qquad s_n=\sqrt{s_n^2}.
\end{equation}
这里分母为\(n\)，是在描述已有数据的平方偏差平均。后面估计未知总体方差时，还会出现分母为\(n-1\)的另一种统计量。

本节数据有\(\bar x=10.5\)、\(s_n^2=4.25\)。若每个值都变为\(y_i=ax_i+b\)，则
\[
 \bar y=a\bar x+b,\qquad s_{y,n}^2=a^2s_{x,n}^2.
\]
还可以用{\heiti 极差}\(x_{(n)}-x_{(1)}\)和{\heiti 四分位距}\(Q_3-Q_1\)描述分散程度。本节二者分别为7和2.5；前者使用两端数据，后者集中于中间一半的位置范围。

\subsection{分组汇总}

设数据分成\(m\)组，第\(h\)组有\(n_h\)个数据，均值为\(\bar x_h\)，描述方差为\(s_h^2\)，总数为\(n=\sum_hn_h\)。直接按组求和，得到
\[
 \bar x=\sum_h\frac{n_h}n\bar x_h.
\]
再写\(x_{hj}-\bar x=(x_{hj}-\bar x_h)+(\bar x_h-\bar x)\)，平方后在各组求和。因为\(\sum_j(x_{hj}-\bar x_h)=0\)，交叉项消失，故
\begin{equation}\label{eq:prob-group-variance}
 s_n^2=\sum_h\frac{n_h}n s_h^2
 +\sum_h\frac{n_h}n(\bar x_h-\bar x)^2.
\end{equation}
合并后的离散程度包含组内波动与组间均值差异。只把各组方差平均，会遗漏第二部分。这与前面条件平均的总方差分解具有相同的有限求和依据。

\begin{example}
甲组20个数据，均值为60、方差为4；乙组30个数据，均值为70、方差为9。求合并后的均值与描述方差，再讨论两组均值相同时公式怎样简化。
\end{example}

\subsection{统计摘要的选择}

仍看式\eqref{eq:prob-description-data}。把最大值15改为30，均值明显上升，而中位数保持不变；极差与平方偏差也对这一变化敏感。若研究总体总量或平方误差，均值和方差很有用；若关注典型位置且存在少量很大的数值，中位数和分位数可以提供另一类信息。

若数据来自不同组别，还需要观察组内分布及其混合比例。同一均值可能来自所有数据集中在中心附近，也可能来自相距很远的两组数据。图形、均值、分位数与离散程度分别保留不同信息，选择摘要时应保留样本量、变量的含义以及数据来源。

\begin{example}[据2025年上海卷改编]\label{ex:prob-swimming-data}
历届奥运会男子\(4\times100\)米混合泳接力冠军的十个成绩数据如下，单位为秒，已按升序排列：
\[
\begin{array}{rrrrr}
206.78&207.46&207.95&209.34&209.35\\
210.68&213.73&214.84&216.93&216.93
\end{array}
\]
求极差、中位数与样本均值；从这十条记录中随机选三条，求恰有两条成绩大于211秒的概率。
\end{example}

\begin{exercise}
\mbox{}\par\nobreak

（1）对数据\(2,3,3,4,5,7,7,9\)，求均值、中位数、众数、四分位距与描述方差，并画箱线图。

（2）（据2020年天津卷第4题改编）由图\ref{fig:prob-histograms}右图，求80个零件中直径在\([5.43,5.47)\)毫米内的个数，并估计该批产品中此类零件的比例。

（3）两组人数相等，均值为\(\mu\pm d\)，方差均为\(v\)。求合并后的均值与方差。

（4）分别构造均值相同而方差不同，以及均值方差相同而形态不同的两组数据。
\end{exercise}
